\section{Rearranged Forward Passes}
\label{app:forward}

A factorised expert layer is useful only if it is never materialised. Written
naively, a product-key layer sums over $H n^2$ terms per token; written
correctly, the same quantity costs $O(H n)$. The rearrangement is where
implementation errors hide, because a wrong contraction still trains and still
drives the loss down --- it simply computes a different function than the one
described. This appendix records the derivations and the check that was run
against them.

Throughout, $\gvec^{(1)}_{hi}$ and $\gvec^{(2)}_{hj}$ are the routing weights
of Equation~\eqref{eq:product-key}, $x \in \R^{d_{\mathrm{in}}}$ is the input,
and indices $i,j$ run over the $n$ factors per side.

\subsection{Horizontal decomposition}

With $E_{ij}$ as in Equation~\eqref{eq:hd}, the layer output is
\begin{equation}
  y = \sum_{h} \sum_{i} \sum_{j}
      \gvec^{(1)}_{hi}\gvec^{(2)}_{hj}
      \Big( V_j\, \sigma(U_i x + b^{(1)}_i) + b^{(2)}_j \Big).
\end{equation}
Write $z_i = \sigma(U_i x + b^{(1)}_i) \in \R^{m}$, computed once for each of
the $n$ bottom factors. The first term factors as
\begin{equation}
  \sum_{h}\sum_{j} \gvec^{(2)}_{hj} V_j
  \underbrace{\sum_{i} \gvec^{(1)}_{hi} z_i}_{t_h \in \R^{m}}
  \;=\;
  \sum_{j} V_j \underbrace{\sum_{h} \gvec^{(2)}_{hj} t_h}_{u_j \in \R^{m}},
\end{equation}
so the bottom factors are collapsed per head into $t_h$ and redistributed to
the top factors as $u_j$, and the $n^2$ expert combinations never appear. The
bias term collapses to $\sum_j c_j b^{(2)}_j$ with $c_j = \sum_h
\gvec^{(2)}_{hj}\sum_i \gvec^{(1)}_{hi}$; when the routing weights are
normalised the inner sum is one, recovering the form given by
\citet{park2025monet}. We implement the general form, which is correct whether
or not normalisation holds.

\subsection{Vertical decomposition}

Here expert $(i,j)$ produces the first half of the output from factor $i$ and
the second half from factor $j$, each acting on both halves of the hidden
state. With $p^{(1)}_i = \sigma(U^{(1)}_i x + b^{(1)}_i)$ and $p^{(2)}_j =
\sigma(U^{(2)}_j x + b^{(2)}_j)$, the two output blocks are
\begin{align}
  y_{\mathrm{top}} &= \sum_i A_i V^{11}_i p^{(1)}_i
                    + \sum_i V^{12}_i \sum_h \gvec^{(1)}_{hi}
                      \sum_j \gvec^{(2)}_{hj} p^{(2)}_j
                    + \sum_i A_i b^{12}_i, \\
  y_{\mathrm{bot}} &= \sum_j B_j V^{22}_j p^{(2)}_j
                    + \sum_j V^{21}_j \sum_h \gvec^{(2)}_{hj}
                      \sum_i \gvec^{(1)}_{hi} p^{(1)}_i
                    + \sum_j B_j b^{22}_j,
\end{align}
where $A_i = \sum_h \gvec^{(1)}_{hi}\sum_j \gvec^{(2)}_{hj}$ and $B_j = \sum_h
\gvec^{(2)}_{hj}\sum_i \gvec^{(1)}_{hi}$. The cross terms are the reason this
variant is more expensive than the horizontal one despite an identical
parameter count: each requires an extra pass through the head index. Since
$P = 113$ is odd, the output is split into blocks of $57$ and $56$.

\subsection{Multilinear layers}

For the CP layer with factors $G^{\mathrm{out}}, G^{\mathrm{in}},
G^{(a)}, G^{(b)}$ of rank $R$,
\begin{equation}
  y = \sum_{h}\sum_{r}
      G^{\mathrm{out}}_{r}
      \big\langle G^{\mathrm{in}}_{r}, x \big\rangle
      \big\langle G^{(a)}_{r}, \gvec^{(1)}_{h} \big\rangle
      \big\langle G^{(b)}_{r}, \gvec^{(2)}_{h} \big\rangle,
\end{equation}
that is, three independent projections multiplied elementwise in rank space.
Tucker replaces the elementwise product by a contraction against a dense core
$\mathcal{Z}$, and Tensor-Train by a chain of three matrix products closed by
$G^{(1)}$. In all three the expert modes are contracted directly against the
routing vectors, so no expert is ever formed. Note that these layers apply no
elementwise nonlinearity between the factors: the output is multilinear in
$(x, \gvec^{(1)}, \gvec^{(2)})$, in contrast to the MONET variants.

\subsection{Numerical verification}

Each rearranged forward pass was re-implemented in NumPy using the same
contraction expressions as the training code, and compared against a naive
implementation that loops over all $H n^2$ expert combinations explicitly.
Table~\ref{tab:verify} reports the maximum absolute deviation on random inputs
and weights. All six agree to machine precision.

\begin{table}[h]
\centering
\caption{Maximum absolute difference between the rearranged forward pass and
the naive per-expert sum, over random inputs and weights.}
\label{tab:verify}
\begin{tabular}{lr}
\toprule
Layer & $\max |y_{\mathrm{fast}} - y_{\mathrm{naive}}|$ \\
\midrule
MONET-HD        & $2.2 \times 10^{-15}$ \\
MONET-VD        & $3.6 \times 10^{-15}$ \\
Unfactorised    & $8.9 \times 10^{-16}$ \\
CP              & $1.1 \times 10^{-14}$ \\
Tucker          & $2.0 \times 10^{-13}$ \\
Tensor-Train    & $4.3 \times 10^{-14}$ \\
\bottomrule
\end{tabular}
\end{table}

\section{Implementation Details}
\label{app:impl}

\paragraph{Analysis without the model.} All spectral analysis is performed on
NumPy from archived weights. The router is re-derived analytically: under
one-hot inputs the logit of key $i$ in head $h$ for operand value $a$ is
exactly the entry $w^{(1)}_{hia}$ of the key matrix, so no forward pass is
needed and the reconstruction is exact rather than sampled. For the
interleaved and random router splits the stored permutation is applied first.

\paragraph{What counts as the shared basis.} For the dense MLP it is the first
weight matrix; for MONET-HD the bottom factor $U$ reshaped to $(nm) \times
d_{\mathrm{in}}$; for MONET-VD the concatenation of $U^{(1)}$ and $U^{(2)}$;
for the multilinear layers the input factor. In each case the block of columns
corresponding to operand $a$ is extracted and Fourier-transformed along that
axis.

\paragraph{Failed runs.} A run that does not reach $95\%$ test accuracy within
$30\,000$ steps is recorded as a failure and excluded from averages over the
grokking step, rather than being assigned the budget as a censored value.
Assigning the budget would understate the gap between architectures that
generalise late and architectures that do not generalise at all.

\paragraph{Numerical conventions.} Spectra are computed after mean removal,
over frequencies $k = 1, \ldots, 56$; the DC component is dropped. Power
spectra are normalised to sum to one per unit before pooling. The white-noise
purity reference of $0.082$ is the mean purity of $4096$ independent Gaussian
vectors of length $113$ under the same procedure.

\paragraph{Compute.} All runs were performed on a single NVIDIA T4. A typical
run takes four to eleven minutes for $30\,000$ full-batch steps on $10\,215$
training examples. Wall-clock time varies little with expert count, since at
this scale the computation is dominated by kernel-launch overhead rather than
arithmetic; observed runtimes therefore do not reflect the asymptotic cost of
the architectures, and the comparison of efficiency in
Table~\ref{tab:params} is made in parameters rather than seconds.